% =====================================================================
%  Mémoire de Master — Détection de la fibrillation auriculaire
%  à partir des intervalles RR par un réseau hybride CNN-LSTM.
%
%  Compilation : pdflatex + biber + pdflatex x2   (ou bouton "Recompile" sur Overleaf)
%  Classe      : report (12pt, A4)
%  Format      : IMRAD strict — directives Sayadi (12pt Times, 1.5, marges 25 mm).
% =====================================================================

\documentclass[12pt,a4paper,oneside]{report}

% ----- Préambule (packages, style, commandes) --------------------------
% =====================================================================
%  Packages
% =====================================================================

% ----- Langue & encodage ---------------------------------------------
\usepackage[utf8]{inputenc}
\usepackage[T1]{fontenc}
\usepackage[french]{babel}
\usepackage{csquotes}                 % requis par biblatex avec babel

% ----- Police Times-like (substitut libre de Times New Roman) --------
\usepackage{newtxtext}                % corps de texte
\usepackage{newtxmath}                % équations cohérentes avec Times

% ----- Mise en page --------------------------------------------------
\usepackage[a4paper,margin=25mm]{geometry}   % marges 25 mm (Sayadi)
\usepackage{setspace}                         % interligne 1.5
\usepackage{indentfirst}                      % premier paragraphe indenté
\usepackage{parskip}                          % petite espace entre paragraphes

% ----- Graphiques ----------------------------------------------------
\usepackage{graphicx}
\graphicspath{{figures/}}
\usepackage{tikz}                            % placeholders figures
\usetikzlibrary{positioning,arrows.meta,shapes.geometric,calc}
\usepackage{float}                            % option [H] pour figures/tableaux

% ----- Tableaux ------------------------------------------------------
\usepackage{booktabs}                         % règles propres : \toprule \midrule
\usepackage{array}                            % colonnes avancées
\usepackage{tabularx}                         % colonnes auto-ajustées
\usepackage{multirow}
\usepackage{caption}
\captionsetup[figure]{labelfont=bf,font=small}
\captionsetup[table]{labelfont=bf,font=small,position=above}

% ----- Mathématiques -------------------------------------------------
\usepackage{amsmath,amsfonts}
\let\Bbbk\relax                               % évite le conflit avec newtxmath
\usepackage{amssymb}
\usepackage{mathtools}
\usepackage{siunitx}                          % unités cohérentes (\SI{0.95}{})

% ----- Algorithmes (encadrés avec titre, sans \fbox) ------------------
\usepackage{algorithm}
\usepackage{algpseudocode}

% ----- Hyperliens ----------------------------------------------------
\usepackage[hidelinks,breaklinks=true]{hyperref}
\usepackage{url}

% ----- Bibliographie (style IEEE, ordre de citation) -----------------
\usepackage[
  backend=biber,
  style=ieee,
  sorting=none,
  doi=false,isbn=false,url=false
]{biblatex}

% ----- En-tête / pied de page ----------------------------------------
\usepackage{fancyhdr}

% ----- Listes & divers -----------------------------------------------
\usepackage{enumitem}
\usepackage{xcolor}
\usepackage{lipsum}                           % texte fictif — à retirer en production

% =====================================================================
%  Style — applique les directives Sayadi
% =====================================================================

% ----- Interligne 1.5 ------------------------------------------------
\onehalfspacing

% ----- En-tête / pied de page ----------------------------------------
\pagestyle{fancy}
\fancyhf{}
\fancyhead[L]{\nouppercase{\leftmark}}
\fancyhead[R]{\thepage}
\renewcommand{\headrulewidth}{0.4pt}
\renewcommand{\footrulewidth}{0pt}

% Page de premier chapitre — style « plain » sans en-tête
\fancypagestyle{plain}{%
  \fancyhf{}%
  \fancyfoot[C]{\thepage}%
  \renewcommand{\headrulewidth}{0pt}%
}

% ----- Mise en forme des titres de chapitre --------------------------
\usepackage{titlesec}
\titleformat{\chapter}[display]
  {\normalfont\Large\bfseries}
  {\chaptertitlename\ \thechapter}{14pt}{\LARGE}
\titlespacing*{\chapter}{0pt}{0pt}{20pt}

\titleformat{\section}{\normalfont\large\bfseries}{\thesection}{1em}{}
\titleformat{\subsection}{\normalfont\normalsize\bfseries}{\thesubsection}{1em}{}

% ----- Justification du texte ----------------------------------------
% Justifié par défaut dans `report` ; on désactive l'over/underfull warnings cosmétiques
\tolerance=1000
\emergencystretch=1em

% ----- Profondeur du sommaire ----------------------------------------
\setcounter{tocdepth}{3}
\setcounter{secnumdepth}{3}

% =====================================================================
%  Commandes personnalisées
% =====================================================================

% ----- Placeholder figure (boîte TikZ avec libellé) -------------------
% Usage : \figplaceholder{Légende} ou \figplaceholder[hauteur=6cm]{Légende}
\NewDocumentCommand{\figplaceholder}{O{0.8\linewidth} O{6cm} m}{%
  \begin{tikzpicture}
    \draw[dashed,line width=0.6pt,fill=gray!8]
      (0,0) rectangle (#1,#2);
    \node[align=center,text width=0.9\dimexpr#1] at (#1/2,#2/2) {%
      \large\textbf{[FIGURE — à insérer]}\\[2pt]%
      \normalsize\itshape #3%
    };
  \end{tikzpicture}%
}

% ----- Notation cohérente VP/VN/FP/FN (directives Critères PDF) ------
\newcommand{\VP}{\mathrm{VP}}
\newcommand{\VN}{\mathrm{VN}}
\newcommand{\FP}{\mathrm{FP}}
\newcommand{\FN}{\mathrm{FN}}

% ----- Acronymes du projet ------------------------------------------
\newcommand{\FA}{FA}                       % Fibrillation auriculaire
\newcommand{\ECG}{ECG}
\newcommand{\CNN}{CNN}
\newcommand{\LSTM}{LSTM}
\newcommand{\AUROC}{AUROC}
\newcommand{\RR}{RR}

% ----- Encadré pour algorithmes (sans \fbox — directive 22 Sayadi) ----
% On utilise l'environnement `algorithm` qui est encadré par convention
% mais sans cadre dessiné autour du contenu.

% ----- Mot-clé visible (mise en valeur sans gras dans le texte) -------
\newcommand{\keyword}[1]{\emph{#1}}

% ----- Numérotation des équations par chapitre ----------------------
\numberwithin{equation}{chapter}
\numberwithin{figure}{chapter}
\numberwithin{table}{chapter}


% ----- Bibliographie ---------------------------------------------------
\addbibresource{bibliography/references.bib}

\begin{document}

% ----- Pages liminaires ------------------------------------------------
\pagenumbering{roman}

% =====================================================================
%  Page de garde — REMPLIR LES CHAMPS [BRACKETS]
% =====================================================================
\begin{titlepage}
  \centering
  \thispagestyle{empty}

  % ----- Bandeau supérieur : université / école / labo ---------------
  {\large \textbf{[NOM DE L'UNIVERSITÉ]}\par}
  \vspace{0.3cm}
  {\large \textbf{[NOM DE L'ÉCOLE / FACULTÉ]}\par}
  \vspace{0.2cm}
  {\normalsize \textit{[DÉPARTEMENT / SPÉCIALITÉ]}\par}

  \vspace{1.0cm}

  % ----- Logos (école à gauche, université à droite) -----------------
  \begin{minipage}{0.3\textwidth}
    \centering
    % \includegraphics[height=3cm]{logos/ecole.png}
    \figplaceholder[3cm][3cm]{Logo école}
  \end{minipage}
  \hfill
  \begin{minipage}{0.3\textwidth}
    \centering
    % \includegraphics[height=3cm]{logos/universite.png}
    \figplaceholder[3cm][3cm]{Logo université}
  \end{minipage}

  \vspace{1.5cm}

  % ----- Type de mémoire --------------------------------------------
  {\Large \textbf{Mémoire de fin d'études}\par}
  \vspace{0.3cm}
  {\large Présenté en vue de l'obtention du diplôme de\par}
  \vspace{0.2cm}
  {\large \textbf{[Master / Ingénieur] en [SPÉCIALITÉ]}\par}

  \vspace{1.5cm}

  % ----- Titre -------------------------------------------------------
  \rule{\linewidth}{0.6pt}
  \vspace{0.3cm}
  {\LARGE\bfseries
    Détection automatique de la fibrillation auriculaire\\
    à partir des intervalles RR\\
    par un réseau hybride CNN--LSTM\par}
  \vspace{0.3cm}
  \rule{\linewidth}{0.6pt}

  \vspace{1.5cm}

  % ----- Auteur ------------------------------------------------------
  {\large Réalisé par\par}
  \vspace{0.2cm}
  {\Large \textbf{Abdelmalek Abed}\par}

  \vspace{1.2cm}

  % ----- Encadrants & jury -------------------------------------------
  \begin{tabular}{ll}
    \textbf{Encadrant académique :} & [Nom Prénom — Grade] \\[4pt]
    \textbf{Encadrant industriel :}  & [Nom Prénom — Organisme] \\[4pt]
  \end{tabular}

  \vspace{0.6cm}

  \textbf{Membres du jury :}
  \begin{tabular}{ll}
    [Nom Prénom] & — Président \\
    [Nom Prénom] & — Rapporteur \\
    [Nom Prénom] & — Examinateur \\
  \end{tabular}

  \vfill

  % ----- Année universitaire ----------------------------------------
  {\large \textbf{Année universitaire 2025--2026}\par}
\end{titlepage}

\cleardoublepage

% =====================================================================
%  Dédicace (optionnelle — remplir ou supprimer)
% =====================================================================
\chapter*{Dédicace}
\thispagestyle{empty}

\begin{flushright}
  \vspace*{3cm}
  \itshape
  À mes parents,\\
  qui ont rendu ce travail possible.\\[1cm]
  À mes proches,\\
  pour leur soutien inconditionnel.\\[1cm]
  À tous mes enseignants.\\
\end{flushright}

\vfill

\cleardoublepage

% =====================================================================
%  Remerciements
%  Règle (Buttler) : remercier personnes ayant contribué + institutions de financement.
% =====================================================================
\chapter*{Remerciements}
\addcontentsline{toc}{chapter}{Remerciements}

% --- Paragraphe d'ouverture ------------------------------------------
Au terme de ce travail, je tiens à exprimer ma profonde gratitude à toutes les personnes qui ont contribué, de près ou de loin, à la réussite de ce mémoire.

% --- Encadrant académique --------------------------------------------
Je remercie tout particulièrement [\textbf{Nom de l'encadrant académique}] pour la confiance qu'il m'a accordée, pour la qualité de son encadrement, et pour ses conseils tout au long de ce projet.

% --- Encadrant industriel / co-encadrant -----------------------------
J'adresse également mes remerciements à [\textbf{Nom de l'encadrant industriel ou co-encadrant}] pour [préciser : sa disponibilité, le partage de son expertise technique, etc.].

% --- Membres du jury -------------------------------------------------
Je remercie les membres du jury, [\textbf{Noms}], d'avoir accepté d'évaluer ce travail et pour le temps qu'ils ont consacré à sa lecture.

% --- Institutions ----------------------------------------------------
Mes remerciements vont aussi à [\textbf{nom de l'école / du laboratoire}] pour les moyens mis à ma disposition durant la réalisation de ce projet.

% --- Proches ---------------------------------------------------------
Enfin, je remercie ma famille et mes amis pour leur soutien constant.

\vfill

\cleardoublepage

% =====================================================================
%  Résumé (FR) + Abstract (EN) + Mots-clés
%  Règle (Buttler) : <= 250 mots, mini-version du mémoire, écrit au passé,
%  autosuffisant, sans citation ni figure.
% =====================================================================

% ---------- Résumé en français ----------------------------------------
\chapter*{Résumé}
\addcontentsline{toc}{chapter}{Résumé}

La fibrillation auriculaire (FA) est l'arythmie cardiaque la plus fréquente et représente un facteur de risque majeur d'accident vasculaire cérébral. Sa détection automatique à partir d'enregistrements ECG longue durée constitue un enjeu clinique reconnu.

Ce travail a proposé une chaîne complète de détection de la FA fondée uniquement sur les intervalles \RR{}, en évitant le coût de l'analyse morphologique de l'\ECG{}. Quatre méthodes de référence ont d'abord été établies sous validation croisée à l'échelle du patient afin de quantifier la difficulté intrinsèque du problème. Un réseau hybride \CNN--\LSTM{} a ensuite été conçu, optimisé par recherche d'hyperparamètres bayésienne (\textit{Optuna}), puis comprimé en vue d'un déploiement embarqué. Sa robustesse a été évaluée par transfert sur un second jeu de données acquis dans des conditions cliniques différentes (LTAFDB). Une application interactive a finalement été développée pour permettre l'exploration des décisions du modèle.

Les résultats ont montré [\textit{à compléter — métriques finales : F1/patient, AUROC, taille modèle après quantification, perte cross-dataset}]. L'analyse a mis en évidence un \textit{plateau intrinsèque} au signal \RR{} seul, dépassable uniquement par l'apport d'une information morphologique additionnelle.

\medskip
\textbf{Mots-clés :} fibrillation auriculaire ; intervalles \RR{} ; apprentissage profond ; réseau hybride \CNN--\LSTM{} ; compression de modèles ; généralisation cross-dataset.

\bigskip

% ---------- Abstract in English ---------------------------------------
\section*{Abstract}

Atrial fibrillation (AF) is the most common cardiac arrhythmia and a major risk factor for stroke. Automated detection from long-term ECG recordings is a recognised clinical need.

This work proposed a complete AF detection pipeline based solely on \RR{} intervals, avoiding the cost of ECG morphology analysis. Four baseline methods were first established under patient-level cross-validation to quantify the intrinsic difficulty of the task. A hybrid \CNN--\LSTM{} network was then designed, tuned via Bayesian hyperparameter search (\textit{Optuna}), and compressed for embedded deployment. Its robustness was assessed through transfer to a second dataset acquired under different clinical conditions (LTAFDB). An interactive application was finally developed to explore the model's decisions.

The results showed [\textit{to be completed — final metrics: patient-level F1, AUROC, post-quantisation model size, cross-dataset performance drop}]. The analysis highlighted an \textit{intrinsic plateau} of \RR{}-only signals, which can only be overcome by adding morphological information.

\medskip
\textbf{Keywords:} atrial fibrillation; \RR{} intervals; deep learning; hybrid \CNN--\LSTM{}; model compression; cross-dataset generalisation.

\cleardoublepage

\tableofcontents
\cleardoublepage

\listoffigures
\cleardoublepage

\listoftables
\cleardoublepage

% Liste des abréviations (optionnelle — remplir au besoin)
% =====================================================================
%  Liste des abréviations
% =====================================================================
\chapter*{Liste des abréviations}
\addcontentsline{toc}{chapter}{Liste des abréviations}

\begin{tabular}{ll}
  \textbf{AFDB}   & MIT-BIH Atrial Fibrillation Database \\
  \textbf{AUROC}  & Area Under the Receiver Operating Characteristic curve \\
  \textbf{AVC}    & Accident Vasculaire Cérébral \\
  \textbf{CNN}    & Convolutional Neural Network \\
  \textbf{CV}     & Cross-Validation (validation croisée) \\
  \textbf{DICE}   & Dice Similarity Coefficient \\
  \textbf{ECG}    & Électrocardiogramme \\
  \textbf{FA}     & Fibrillation Auriculaire \\
  \textbf{FN}     & Faux Négatif \\
  \textbf{FP}     & Faux Positif \\
  \textbf{INT8}   & Quantification entière sur 8 bits \\
  \textbf{LSTM}   & Long Short-Term Memory \\
  \textbf{LTAFDB} & Long-Term AF Database \\
  \textbf{ONNX}   & Open Neural Network Exchange \\
  \textbf{ROC}    & Receiver Operating Characteristic \\
  \textbf{RR}     & Intervalle entre deux ondes R successives \\
  \textbf{VN}     & Vrai Négatif \\
  \textbf{VP}     & Vrai Positif \\
\end{tabular}

\cleardoublepage

% ----- Corps du mémoire ------------------------------------------------
\pagenumbering{arabic}
\setcounter{page}{1}

% =====================================================================
%  Chapitre 1 — Introduction
%  Règles (Sayadi + Buttler) :
%   - temps : présent
%   - structure entonnoir : contexte → état de l'art bref (chaque travail
%     cité en quelques lignes avec [n]) → problème → objectifs →
%     contributions → plan chapitre par chapitre.
% =====================================================================
\chapter{Introduction}
\label{ch:intro}

% --- 1.1 Contexte général --------------------------------------------
\section{Contexte général}
\label{sec:intro:contexte}

La fibrillation auriculaire (\FA{}) est l'arythmie cardiaque soutenue la plus fréquente. Elle multiplie par cinq le risque d'accident vasculaire cérébral et constitue une cause majeure d'hospitalisation \cite{hindricks2021guidelines}. Son diagnostic repose sur l'électrocardiogramme (\ECG{}), enregistré soit en cabinet, soit en continu à l'aide d'enregistreurs ambulatoires de type \textit{Holter}.

L'analyse manuelle de tels enregistrements, qui peuvent dépasser vingt-quatre heures, demeure chronophage. Les progrès récents en \keyword{apprentissage profond} ont ouvert la voie à des outils d'aide à la décision capables de signaler automatiquement les épisodes de \FA{}. Parmi les caractéristiques utilisables, les \keyword{intervalles \RR{}} — durées entre deux battements consécutifs — présentent un intérêt particulier : ils se calculent à partir d'un seul point de référence (l'onde R) et restent robustes au bruit et aux variations d'amplitude du signal \ECG{}.

% --- 1.2 État de l'art (synthétique) ---------------------------------
\section{État de l'art en bref}
\label{sec:intro:soa}

Plusieurs familles de méthodes ont été proposées dans la littérature.

Les approches classiques reposent sur des descripteurs statistiques de la série \RR{} suivis d'un classifieur supervisé : \citeauthor{tateno2001automatic} exploitent un test de coefficient de variation et la matrice de transition des \RR{} \cite{tateno2001automatic} ; \citeauthor{lake2011accurate} utilisent une entropie multi-échelle pour la détection \cite{lake2011accurate}.

Les approches profondes récentes traitent directement le signal \ECG{} brut. \citeauthor{hannun2019cardiologist} ont entraîné un \CNN{} à 34 couches sur plus de 90~000 \ECG{} ambulatoires et atteint un \AUROC{} de 0{,}97 \cite{hannun2019cardiologist}. Sur le seul signal \RR{}, \citeauthor{andersen2019deep} proposent une architecture hybride \CNN--\LSTM{} et rapportent une sensibilité supérieure à 98~\% en classification de fenêtre \cite{andersen2019deep}.

La question de la \keyword{robustesse cross-dataset} — c'est-à-dire le maintien des performances d'un modèle entraîné sur un jeu de données et évalué sur un autre — est, en revanche, peu étudiée. De même, la contrainte de \keyword{compacité} requise pour un déploiement embarqué n'est pas systématiquement examinée.

% --- 1.3 Problématique & objectifs -----------------------------------
\section{Problématique et objectifs}
\label{sec:intro:objectifs}

Ce travail s'intéresse à la détection automatique de la \FA{} \textbf{à partir des seuls intervalles \RR{}}, en respectant trois exigences pratiques :
\begin{enumerate}[leftmargin=*]
  \item \textbf{Évaluation honnête} : isoler la difficulté intrinsèque du problème par une validation croisée à l'échelle du \emph{patient}, et non de la fenêtre.
  \item \textbf{Compacité} : produire un modèle final dont la taille permet un déploiement sur dispositif portable.
  \item \textbf{Robustesse} : quantifier la perte de performance lors d'un transfert vers un jeu de données acquis dans des conditions cliniques différentes.
\end{enumerate}

% --- 1.4 Contributions -----------------------------------------------
\section{Contributions}
\label{sec:intro:contributions}

Ce mémoire apporte les contributions suivantes :
\begin{itemize}[leftmargin=*]
  \item Une étude comparative de quatre familles de classifieurs (règle métier, gradient boosting, \CNN{} 1D, \CNN--\LSTM{}) sous protocole patient-level identique.
  \item Une architecture \CNN--\LSTM{} optimisée par \textit{Optuna} avec analyse d'ablation systématique.
  \item Une étude de compression couvrant huit variantes (pruning, quantification dynamique et INT8, export ONNX) avec analyse des compromis taille / latence / performance.
  \item Une évaluation cross-dataset \textbf{AFDB~$\rightarrow$~LTAFDB} mettant en évidence un plateau intrinsèque au signal \RR{} seul.
  \item Un tableau de bord interactif \textit{Streamlit} permettant l'exploration patient par patient des décisions du modèle.
\end{itemize}

% --- 1.5 Plan du mémoire (obligatoire — Sayadi 1) --------------------
\section{Plan du mémoire}
\label{sec:intro:plan}

Le présent document est organisé en six chapitres.

Le \textbf{chapitre~\ref{ch:soa}} présente l'état de l'art en détection automatique de la \FA{} : éléments cliniques, jeux de données publics et approches algorithmiques.

Le \textbf{chapitre~\ref{ch:methodes}} décrit la méthodologie adoptée : pipeline de prétraitement, architecture proposée, protocole d'évaluation à l'échelle du patient et critères statistiques (\VP{}, \VN{}, \FP{}, \FN{}, précision, rappel, F1, courbe ROC).

Le \textbf{chapitre~\ref{ch:resultats}} rapporte les résultats expérimentaux des six phases du projet : exploration des données, baselines, modèle \CNN--\LSTM{}, étude de compression, robustesse cross-dataset et application interactive.

Le \textbf{chapitre~\ref{ch:discussion}} discute ces résultats, les compare à la littérature et identifie les limites du travail.

Le \textbf{chapitre~\ref{ch:conclusion}} conclut le mémoire et propose plusieurs perspectives.

% =====================================================================
%  Chapitre 2 — État de l'art
%  Règles (Sayadi/Buttler) :
%   - structure : du général au particulier, transitions entre parties
%   - chaque travail cité = quelques lignes + [n]
%   - temps : présent (faits établis), passé (allusions aux travaux)
% =====================================================================
\chapter{État de l'art}
\label{ch:soa}

% --- 2.1 Rappels cliniques -------------------------------------------
\section{Rappels cliniques sur la fibrillation auriculaire}
\label{sec:soa:clinique}

\subsection{Définition et impact}
La \FA{} se caractérise par une activité électrique désorganisée des oreillettes, entraînant un rythme ventriculaire irrégulier. Sa prévalence augmente avec l'âge et atteint près de 10~\% au-delà de 80~ans \cite{hindricks2021guidelines}.

\subsection{Signature électrocardiographique}
Sur l'\ECG{}, la \FA{} se traduit par trois signes principaux :
\begin{itemize}
  \item l'\textbf{absence d'onde P} ;
  \item la présence d'oscillations dites \textbf{de fibrillation} ;
  \item une \textbf{irrégularité des intervalles \RR{}}.
\end{itemize}

\begin{figure}[H]
  \centering
  \figplaceholder[0.85\linewidth][6cm]{ECG en rythme sinusal vs ECG en FA — illustration des trois signes}
  \caption{Exemples d'\ECG{} en rythme sinusal (haut) et en fibrillation auriculaire (bas). L'absence d'onde P et l'irrégularité des intervalles \RR{} sont caractéristiques de la \FA{}.}
  \label{fig:soa:ecg_fa}
\end{figure}

% --- 2.2 Signal RR ---------------------------------------------------
\section{Le signal \RR{} comme entrée monodimensionnelle}
\label{sec:soa:rr}

\subsection{Définition}
L'intervalle \RR{} désigne la durée séparant deux ondes R consécutives, c'est-à-dire deux contractions ventriculaires. La suite des intervalles \RR{} forme un \keyword{tachogramme} qui résume l'information de rythme indépendamment de la morphologie de l'\ECG{}.

\subsection{Avantages et limites}
L'analyse fondée sur les \RR{} présente plusieurs avantages :
\begin{itemize}
  \item robustesse au bruit basse fréquence et aux variations d'amplitude ;
  \item faible coût computationnel ;
  \item compatibilité avec les capteurs portables.
\end{itemize}
Elle souffre toutefois d'une limite intrinsèque : elle ignore la morphologie de l'onde P, dont l'absence est pourtant un marqueur direct de la \FA{}.

% --- 2.3 Jeux de données ---------------------------------------------
\section{Jeux de données publics}
\label{sec:soa:datasets}

\subsection{MIT-BIH Atrial Fibrillation Database (AFDB)}
La base AFDB \cite{moody1983new} contient 25~enregistrements \ECG{} ambulatoires de 10~heures, échantillonnés à 250~Hz, annotés battement par battement. Elle est devenue le \textit{benchmark} de référence pour la détection de la \FA{}.

\subsection{Long-Term AF Database (LTAFDB)}
La base LTAFDB \cite{petrutiu2007abrupt} contient 84~enregistrements de 24 à 25~heures, acquis dans des conditions cliniques différentes. Elle est utilisée dans ce travail comme jeu de test \emph{externe} pour évaluer la robustesse cross-dataset.

\begin{table}[H]
  \centering
  \caption{Caractéristiques des deux bases utilisées dans ce mémoire.}
  \label{tab:soa:datasets}
  \begin{tabular}{lcc}
    \toprule
    & \textbf{AFDB} & \textbf{LTAFDB} \\
    \midrule
    Nombre de patients          & 25 & 84 \\
    Durée par enregistrement    & 10 h & 24--25 h \\
    Fréquence d'échantillonnage & 250 Hz & 128 Hz \\
    Annotations                 & battement & battement \\
    Usage dans ce mémoire       & train / val / test & test externe \\
    \bottomrule
  \end{tabular}
\end{table}

% --- 2.4 Approches algorithmiques ------------------------------------
\section{Approches algorithmiques pour la détection de la \FA{}}
\label{sec:soa:approches}

\subsection{Méthodes statistiques}
Les premières approches ont exploité des descripteurs statistiques de la série \RR{} : coefficient de variation, racine de la moyenne des carrés des différences successives (RMSSD), entropie de Shannon, entropie multi-échelle \cite{tateno2001automatic,lake2011accurate}.

\subsection{Apprentissage automatique classique}
Avec l'essor de l'apprentissage automatique, ces descripteurs ont alimenté des classifieurs supervisés tels que les forêts aléatoires et le gradient boosting \cite{XX_baseline_gb}.

\subsection{Apprentissage profond sur \ECG{} brut}
Les architectures profondes appliquées directement au signal \ECG{} ont récemment atteint une performance comparable à celle d'un cardiologue \cite{hannun2019cardiologist}.

\subsection{Apprentissage profond sur intervalles \RR{}}
Plus pertinents pour ce mémoire, les travaux ciblant le seul signal \RR{} reposent généralement sur des architectures hybrides combinant convolutions 1D et couches récurrentes \cite{andersen2019deep}.

% --- 2.5 Synthèse & positionnement -----------------------------------
\section{Synthèse et positionnement}
\label{sec:soa:positionnement}

La littérature révèle trois éléments insuffisamment traités :
\begin{enumerate}[leftmargin=*]
  \item l'\textbf{évaluation à l'échelle du patient}, souvent remplacée par une évaluation par fenêtre, qui surestime la performance ;
  \item l'\textbf{empreinte mémoire} du modèle final, rarement rapportée ;
  \item la \textbf{généralisation cross-dataset}, peu étudiée.
\end{enumerate}
Le présent travail se positionne précisément sur ces trois axes.

% =====================================================================
%  Chapitre 3 — Matériel et méthodes
%  Règles (Sayadi/Buttler) :
%   - temps : passé
%   - algorithmes séparés, encadrés, avec titre (env. `algorithm`)
%   - organigrammes/schémas plutôt que phrases pour les pipelines
%   - reproductibilité totale
% =====================================================================
\chapter{Matériel et méthodes}
\label{ch:methodes}

% --- 3.1 Vue d'ensemble ---------------------------------------------
\section{Vue d'ensemble du pipeline}
\label{sec:meth:pipeline}

Le pipeline complet, illustré à la figure~\ref{fig:meth:pipeline}, a été décomposé en six étapes successives.

\begin{figure}[H]
  \centering
  \figplaceholder[0.95\linewidth][7cm]{Organigramme du pipeline : AFDB $\rightarrow$ détection R $\rightarrow$ fenêtrage $\rightarrow$ split patient-level $\rightarrow$ modèle $\rightarrow$ évaluation}
  \caption{Vue d'ensemble du pipeline de détection de la \FA{} adopté dans ce mémoire.}
  \label{fig:meth:pipeline}
\end{figure}

% --- 3.2 Données -----------------------------------------------------
\section{Données et prétraitement}
\label{sec:meth:donnees}

\subsection{Source et accès}
Les enregistrements ont été téléchargés depuis PhysioNet via la bibliothèque \texttt{wfdb}. Pour chaque enregistrement, les annotations de battements et de rythme ont été récupérées au format natif.

\subsection{Construction du tachogramme}
Pour chaque patient, la suite des intervalles \RR{} a été obtenue à partir des positions d'ondes R fournies par les annotations expertes. Aucune détection automatique de R n'a été réalisée afin d'éviter d'introduire un biais lié au détecteur.

\subsection{Fenêtrage}
La série \RR{} a été segmentée en fenêtres glissantes de longueur \(W\) battements. Chaque fenêtre a reçu l'étiquette du rythme majoritaire couvert par celle-ci (un seuil de pureté de 50~\% a été appliqué).

\subsection{Split patient-level}
Tous les patients ont été assignés exclusivement à un sous-ensemble (entraînement, validation ou test) : aucune fuite inter-patient n'a été tolérée. Ce protocole, central pour l'honnêteté de l'évaluation, est illustré par l'algorithme~\ref{alg:meth:split}.

\begin{algorithm}[H]
  \caption{Validation croisée à l'échelle du patient}
  \label{alg:meth:split}
  \begin{algorithmic}[1]
    \Require ensemble de patients $\mathcal{P}$, nombre de \textit{folds} $K$
    \Ensure \textit{folds} de patients $(\mathcal{F}_1, \dots, \mathcal{F}_K)$ disjoints
    \State Mélanger $\mathcal{P}$ avec graine fixée
    \State Partitionner $\mathcal{P}$ en $K$ sous-ensembles équilibrés en prévalence FA
    \For{$k = 1$ à $K$}
      \State $\mathcal{F}_k^{\text{test}} \gets$ patients du \textit{fold} $k$
      \State $\mathcal{F}_k^{\text{train}} \gets \mathcal{P} \setminus \mathcal{F}_k^{\text{test}}$
    \EndFor
    \State \Return $(\mathcal{F}_1, \dots, \mathcal{F}_K)$
  \end{algorithmic}
\end{algorithm}

% --- 3.3 Modèles -----------------------------------------------------
\section{Modèles étudiés}
\label{sec:meth:modeles}

\subsection{Baselines}
Quatre méthodes de référence ont été implémentées :
\begin{itemize}
  \item \textbf{Règle métier} fondée sur le coefficient de variation des \RR{} ;
  \item \textbf{Forêt aléatoire} sur descripteurs statistiques ;
  \item \textbf{Gradient boosting} (XGBoost) sur les mêmes descripteurs ;
  \item \textbf{\CNN{} 1D} appliqué directement à la série \RR{}.
\end{itemize}

\subsection{Modèle proposé : \CNN--\LSTM{}}
Le modèle final combine un encodeur convolutionnel 1D et une couche récurrente \LSTM{} bidirectionnelle. Son architecture est décrite à la figure~\ref{fig:meth:archi}.

\begin{figure}[H]
  \centering
  \figplaceholder[0.85\linewidth][7cm]{Schéma de l'architecture CNN-LSTM (encodeur conv \(\rightarrow\) BiLSTM \(\rightarrow\) tête dense)}
  \caption{Architecture du modèle hybride \CNN--\LSTM{} proposé.}
  \label{fig:meth:archi}
\end{figure}

\subsection{Recherche d'hyperparamètres}
Les hyperparamètres ont été optimisés à l'aide de la bibliothèque \texttt{Optuna} \cite{akiba2019optuna}. La fonction objectif a été définie comme le \(F_1\) moyenné sur les \textit{folds} de validation.

\subsection{Compression}
Quatre techniques ont été combinées : pruning structuré, quantification dynamique PyTorch, quantification statique INT8 et export ONNX. Les variantes étudiées sont synthétisées à la table~\ref{tab:meth:compression}.

\begin{table}[H]
  \centering
  \caption{Variantes de compression évaluées.}
  \label{tab:meth:compression}
  \begin{tabular}{lll}
    \toprule
    \textbf{Variante} & \textbf{Technique} & \textbf{Cible} \\
    \midrule
    V0 & Modèle de référence FP32 & Référence \\
    V1 & Pruning structuré 30~\% & Sparsité \\
    V2 & Pruning + fine-tune & Récupération performance \\
    V3 & Quantification dynamique & Inférence CPU \\
    V4 & Quantification statique INT8 & Empreinte mémoire \\
    V5 & Export ONNX FP32 & Portabilité \\
    V6 & ONNX + quantification & Embarqué \\
    V7 & Combinaison pruning + INT8 & Taille minimale \\
    \bottomrule
  \end{tabular}
\end{table}

% --- 3.4 Critères d'évaluation -- IMPORTANT : intégrer le PDF Critères --
\section{Critères d'évaluation}
\label{sec:meth:criteres}

L'évaluation d'un classifieur binaire repose sur la matrice de confusion présentée à la table~\ref{tab:meth:confusion}.

\begin{table}[H]
  \centering
  \caption{Matrice de confusion d'un classifieur binaire.}
  \label{tab:meth:confusion}
  \begin{tabular}{c|cc}
    \toprule
     & \textbf{Prédit négatif} & \textbf{Prédit positif} \\
    \midrule
    \textbf{Réel négatif} & \VN{} & \FP{} \\
    \textbf{Réel positif} & \FN{} & \VP{} \\
    \bottomrule
  \end{tabular}
\end{table}

\noindent où :
\begin{itemize}[leftmargin=*]
  \item \(\VP{}\) : nombre de cas correctement classés positifs ;
  \item \(\VN{}\) : nombre de cas correctement classés négatifs ;
  \item \(\FP{}\) : nombre de cas négatifs incorrectement classés positifs ;
  \item \(\FN{}\) : nombre de cas positifs incorrectement classés négatifs.
\end{itemize}

Dans le contexte de la détection de la \FA{}, la minimisation des \FN{} est prioritaire : un épisode manqué peut conduire à un accident vasculaire cérébral non prévenu.

\subsection{Métriques classiques}

\begin{equation}
  \mathrm{Exactitude} = \frac{\VP{} + \VN{}}{\VP{} + \VN{} + \FP{} + \FN{}}
  \label{eq:meth:accuracy}
\end{equation}

\begin{equation}
  \mathrm{Pr\acute{e}cision} = \frac{\VP{}}{\VP{} + \FP{}}
  \label{eq:meth:precision}
\end{equation}

\begin{equation}
  \mathrm{Rappel} = \mathrm{Sensibilit\acute{e}} = \frac{\VP{}}{\VP{} + \FN{}}
  \label{eq:meth:rappel}
\end{equation}

\begin{equation}
  \mathrm{Sp\acute{e}cificit\acute{e}} = \frac{\VN{}}{\VN{} + \FP{}}
  \label{eq:meth:specificite}
\end{equation}

\begin{equation}
  F_1 = 2 \cdot \frac{\mathrm{Pr\acute{e}cision} \cdot \mathrm{Rappel}}{\mathrm{Pr\acute{e}cision} + \mathrm{Rappel}}
  \label{eq:meth:f1}
\end{equation}

\begin{equation}
  \mathrm{DICE} = \frac{2\,\VP{}}{2\,\VP{} + \FP{} + \FN{}}
  \label{eq:meth:dice}
\end{equation}

L'exactitude (équation~\ref{eq:meth:accuracy}) perd de sa pertinence en présence de classes déséquilibrées. Dans ce travail, où la prévalence de la \FA{} est inférieure à 50~\%, le score \(F_1\) (équation~\ref{eq:meth:f1}) a été préféré comme critère principal.

\subsection{Courbe ROC et aire sous la courbe (AUC)}
La courbe ROC trace la sensibilité en fonction du complément à 1 de la spécificité, paramétrée par le seuil de décision. L'aire sous la courbe (\AUROC{}) constitue une métrique de séparabilité indépendante du seuil.

\begin{figure}[H]
  \centering
  \figplaceholder[0.7\linewidth][6cm]{Courbe ROC type avec AUC = 0,90 — diagonale = classifieur aléatoire}
  \caption{Exemple de courbe ROC. Plus l'aire sous la courbe est proche de 1, meilleur est le classifieur.}
  \label{fig:meth:roc}
\end{figure}

\subsection{Représentation par boîte à moustaches}
La performance étant agrégée à l'échelle du \emph{patient}, la distribution des scores inter-patients a été visualisée par des boîtes à moustaches afin de mettre en évidence la dispersion plutôt que la seule moyenne.

% --- 3.5 Environnement -----------------------------------------------
\section{Environnement logiciel et matériel}
\label{sec:meth:env}

Les expériences ont été menées en Python~3.11 avec PyTorch~2.x, scikit-learn et Optuna. Le matériel utilisé est résumé à la table~\ref{tab:meth:env} ; le fichier \texttt{requirements-lock.txt} du dépôt fixe les versions exactes pour assurer la reproductibilité.

\begin{table}[H]
  \centering
  \caption{Environnement logiciel et matériel.}
  \label{tab:meth:env}
  \begin{tabular}{ll}
    \toprule
    Système d'exploitation & Linux 6.8 \\
    Langage                & Python 3.11 \\
    Frameworks             & PyTorch 2.x, scikit-learn, Optuna, ONNX Runtime \\
    GPU                    & [À PRÉCISER] \\
    CPU                    & [À PRÉCISER] \\
    Suivi d'expériences    & Optuna SQLite \\
    \bottomrule
  \end{tabular}
\end{table}

% =====================================================================
%  Chapitre 4 — Résultats
%  Règles (Sayadi 4) :
%   - tableaux/courbes/figures totalement DÉCRITS (variations, valeurs, taux d'erreur)
%   - et COMMENTÉS (comparaison, amélioration, limitation)
%   - éviter de dupliquer figure + tableau
%   - temps : passé pour l'analyse, présent pour décrire une figure/tableau
% =====================================================================
\chapter{Résultats expérimentaux}
\label{ch:resultats}

% --- 4.1 Phase 1 — EDA ----------------------------------------------
\section{Phase 1 — Exploration des données}
\label{sec:res:eda}

\subsection{Prévalence de la \FA{} par patient}
La répartition de la prévalence de la \FA{} parmi les 25~patients d'AFDB est rapportée à la figure~\ref{fig:res:prevalence}. La prévalence par patient s'étend de \texttt{[XX]\%} à \texttt{[XX]\%} avec une médiane de \texttt{[XX]\%}, ce qui confirme une \keyword{forte hétérogénéité inter-patients}.

\begin{figure}[H]
  \centering
  \figplaceholder[0.85\linewidth][6cm]{Histogramme + boîte à moustaches de la prévalence FA par patient (AFDB)}
  \caption{Distribution de la prévalence de la \FA{} par patient dans la base AFDB.}
  \label{fig:res:prevalence}
\end{figure}

\subsection{Distribution des intervalles \RR{}}
Les distributions conditionnelles des \RR{} en rythme sinusal et en \FA{} sont présentées à la figure~\ref{fig:res:rr_dist}. La distribution en \FA{} est sensiblement plus étalée et plus asymétrique, ce qui justifie l'utilisation d'un descripteur de variabilité comme entrée du modèle.

\begin{figure}[H]
  \centering
  \figplaceholder[0.85\linewidth][6cm]{Densités RR : rythme sinusal vs FA (kernel density estimate)}
  \caption{Densités des intervalles \RR{} conditionnellement au rythme.}
  \label{fig:res:rr_dist}
\end{figure}

% --- 4.2 Phase 2 — Baselines -----------------------------------------
\section{Phase 2 — Méthodes de référence}
\label{sec:res:baselines}

Les performances des quatre baselines sous validation croisée à 5~plis patient-level sont rapportées à la table~\ref{tab:res:baselines}.

\begin{table}[H]
  \centering
  \caption{Performance des baselines (médiane $\pm$ IQR sur 25~patients, AFDB).}
  \label{tab:res:baselines}
  \begin{tabular}{lcccc}
    \toprule
    \textbf{Méthode} & \textbf{F1 patient} & \textbf{\AUROC{}} & \textbf{Sensibilité} & \textbf{Spécificité} \\
    \midrule
    Règle métier            & [XX] & [XX] & [XX] & [XX] \\
    Forêt aléatoire         & [XX] & [XX] & [XX] & [XX] \\
    Gradient boosting       & [XX] & [XX] & [XX] & [XX] \\
    \CNN{} 1D               & [XX] & [XX] & [XX] & [XX] \\
    \bottomrule
  \end{tabular}
\end{table}

Les résultats montrent que \texttt{[À COMPLÉTER]}. Une remarque centrale ressort : \emph{aucune} des baselines ne dépasse un \(F_1\) patient de \texttt{[XX]}, ce qui suggère un plafond intrinsèque difficile à franchir avec ces seules approches.

% --- 4.3 Phase 3 — CNN-LSTM + Optuna --------------------------------
\section{Phase 3 — Modèle \CNN--\LSTM{} optimisé}
\label{sec:res:cnnlstm}

\subsection{Recherche d'hyperparamètres}
Optuna a exploré \texttt{[XX]} essais sur \texttt{[XX]} heures. La distribution des scores et l'importance relative des hyperparamètres sont présentées à la figure~\ref{fig:res:optuna}.

\begin{figure}[H]
  \centering
  \figplaceholder[0.85\linewidth][6cm]{Diagramme Optuna : importance des hyperparamètres + scatter score vs nombre d'essais}
  \caption{Résultats de la recherche bayésienne par \textit{Optuna}.}
  \label{fig:res:optuna}
\end{figure}

\subsection{Performance du modèle retenu}
La table~\ref{tab:res:cnnlstm} compare le modèle retenu aux meilleures baselines.

\begin{table}[H]
  \centering
  \caption{Comparaison du \CNN--\LSTM{} optimisé aux baselines.}
  \label{tab:res:cnnlstm}
  \begin{tabular}{lcccc}
    \toprule
    \textbf{Méthode} & \textbf{F1 patient} & \textbf{\AUROC{}} & \textbf{Sens.} & \textbf{Spéc.} \\
    \midrule
    Meilleure baseline (GB) & [XX] & [XX] & [XX] & [XX] \\
    \CNN--\LSTM{} (initial) & [XX] & [XX] & [XX] & [XX] \\
    \CNN--\LSTM{} (Optuna)  & [XX] & [XX] & [XX] & [XX] \\
    \bottomrule
  \end{tabular}
\end{table}

L'amélioration apportée par l'optimisation Optuna est de \texttt{[XX]} points de \(F_1\) en valeur absolue par rapport au modèle initial.

\subsection{Étude d'ablation}
Une étude d'ablation a été menée pour quantifier la contribution de chaque composant ; ses résultats apparaissent à la figure~\ref{fig:res:ablation}.

\begin{figure}[H]
  \centering
  \figplaceholder[0.85\linewidth][6cm]{Barplot ablation : F1 obtenu en retirant chaque composant (régularisation, BiLSTM, dropout, etc.)}
  \caption{Étude d'ablation des composants du modèle.}
  \label{fig:res:ablation}
\end{figure}

% --- 4.4 Phase 4 — Compression --------------------------------------
\section{Phase 4 — Étude de compression}
\label{sec:res:compression}

Les huit variantes de compression évaluées sont synthétisées à la table~\ref{tab:res:compression}.

\begin{table}[H]
  \centering
  \caption{Compromis taille / performance des variantes de compression.}
  \label{tab:res:compression}
  \begin{tabular}{lcccc}
    \toprule
    \textbf{Variante} & \textbf{Taille (KB)} & \textbf{F1 patient} & \(\Delta\)F1 & \textbf{Latence (ms)} \\
    \midrule
    V0 — référence FP32          & [XX] & [XX] & 0 & [XX] \\
    V1 — pruning 30~\%           & [XX] & [XX] & [XX] & [XX] \\
    V2 — pruning + fine-tune     & [XX] & [XX] & [XX] & [XX] \\
    V3 — quant. dynamique        & [XX] & [XX] & [XX] & [XX] \\
    V4 — quant. statique INT8    & 87   & [XX] & [XX] & [XX] \\
    V5 — ONNX FP32               & [XX] & [XX] & [XX] & [XX] \\
    V6 — ONNX + quant.           & [XX] & [XX] & [XX] & [XX] \\
    V7 — pruning + INT8          & [XX] & [XX] & [XX] & [XX] \\
    \bottomrule
  \end{tabular}
\end{table}

La variante V4 a atteint la cible des 200~KB visée pour l'embarqué (87~KB obtenus) avec une perte de \(F_1\) inférieure à \texttt{[XX]} points. Le compromis taille / performance est visualisé à la figure~\ref{fig:res:pareto}.

\begin{figure}[H]
  \centering
  \figplaceholder[0.85\linewidth][6cm]{Frontière Pareto : F1 vs taille pour les 8 variantes}
  \caption{Compromis taille / performance des variantes de compression.}
  \label{fig:res:pareto}
\end{figure}

% --- 4.5 Phase 5 — Cross-dataset ------------------------------------
\section{Phase 5 — Robustesse cross-dataset (AFDB $\rightarrow$ LTAFDB)}
\label{sec:res:cross}

Le modèle entraîné exclusivement sur AFDB a été évalué sans réentraînement sur LTAFDB. Les résultats apparaissent à la table~\ref{tab:res:cross}.

\begin{table}[H]
  \centering
  \caption{Performance du modèle final sur les deux jeux de données.}
  \label{tab:res:cross}
  \begin{tabular}{lccc}
    \toprule
    \textbf{Jeu de test} & \textbf{F1 patient} & \textbf{\AUROC{}} & \textbf{Sens.} \\
    \midrule
    AFDB (interne, CV)        & [XX] & [XX] & [XX] \\
    LTAFDB (externe)          & [XX] & [XX] & [XX] \\
    \(\Delta\) absolu         & [XX] & [XX] & [XX] \\
    \bottomrule
  \end{tabular}
\end{table}

La perte cross-dataset observée est de \texttt{[XX]} points de \(F_1\). La figure~\ref{fig:res:boxplot_cross} montre la dispersion par patient sur les deux bases.

\begin{figure}[H]
  \centering
  \figplaceholder[0.85\linewidth][6cm]{Boîtes à moustaches F1/patient : AFDB vs LTAFDB côte à côte}
  \caption{Dispersion inter-patient sur les deux jeux de données.}
  \label{fig:res:boxplot_cross}
\end{figure}

% --- 4.6 Phase 6 — Sandbox ------------------------------------------
\section{Phase 6 — Application interactive}
\label{sec:res:sandbox}

Une application \textit{Streamlit} a été développée pour permettre l'exploration patient par patient des décisions du modèle. Elle comporte cinq onglets : \emph{Patient Inspector}, \emph{Model Showdown}, \emph{Compression Lab}, \emph{Cross-Dataset} et \emph{What-If}.

\begin{figure}[H]
  \centering
  \figplaceholder[0.85\linewidth][7cm]{Capture d'écran de l'onglet Patient Inspector — tachogramme + probabilité prédite}
  \caption{Capture d'écran de l'application \textit{Streamlit} (\emph{Patient Inspector}).}
  \label{fig:res:sandbox}
\end{figure}

% =====================================================================
%  Chapitre 5 — Discussion
%  Règles (Sayadi/Buttler) :
%   - temps : présent
%   - trier les faits par importance, démonstration
%   - comparer à la littérature
%   - identifier limitations
%   - transitions claires entre paragraphes
% =====================================================================
\chapter{Discussion}
\label{ch:discussion}

% --- 5.1 Synthèse des résultats principaux ---------------------------
\section{Synthèse des résultats principaux}
\label{sec:disc:synthese}

Les trois résultats centraux de ce travail sont les suivants.

\textbf{Premièrement}, sous une évaluation patient-level rigoureuse, toutes les méthodes — qu'elles soient classiques ou profondes — convergent vers un \(F_1\) compris entre \texttt{[0,XX]} et \texttt{[0,XX]}. Ce constat suggère l'existence d'un plafond intrinsèque au signal \RR{} seul.

\textbf{Deuxièmement}, l'architecture \CNN--\LSTM{} optimisée par Optuna n'a apporté qu'un gain limité sur ce plafond, tout en bénéficiant d'une marge importante en compression : la variante INT8 a atteint 87~KB sans dégradation significative.

\textbf{Troisièmement}, le transfert AFDB~$\rightarrow$~LTAFDB s'est traduit par une perte de \texttt{[XX]} points, confirmant que la généralisation cross-dataset reste un point dur.

% --- 5.2 Comparaison à la littérature -------------------------------
\section{Comparaison à la littérature}
\label{sec:disc:litterature}

La performance brute obtenue (\(F_1 \approx \texttt{[XX]}\) sur AFDB) reste inférieure aux valeurs typiquement annoncées dans la littérature, qui dépassent souvent 0,95 \cite{andersen2019deep,hannun2019cardiologist}. Cet écart trouve une explication méthodologique.

La plupart des travaux rapportent une performance \keyword{par fenêtre} après mélange des fenêtres de tous les patients, ce qui crée une fuite : la même série temporelle apparaît à la fois en entraînement et en test. Une évaluation patient-level — adoptée dans ce mémoire — fait disparaître cette fuite et révèle une difficulté qui n'apparaît pas dans une évaluation par fenêtre.

Lorsque l'on bascule l'évaluation \emph{de ce même modèle} en mode fenêtre, on retrouve des chiffres comparables à la littérature ; la différence n'est donc pas due à l'architecture, mais bien au protocole d'évaluation.

% --- 5.3 Étude de compression ---------------------------------------
\section{Compression : un compromis favorable}
\label{sec:disc:compression}

La table~\ref{tab:res:compression} a montré que la quantification statique INT8 réduit la taille du modèle d'un facteur d'environ \texttt{[XX]} sans pénaliser de plus de \texttt{[XX]} points le \(F_1\). Ce résultat ouvre la voie à un déploiement sur dispositif portable.

Trois pièges méthodologiques ont été rencontrés et méritent d'être signalés :
\begin{enumerate}[leftmargin=*]
  \item le \emph{pruning seul ne réduit pas la taille du fichier} tant que les zéros ne sont pas effectivement supprimés du tenseur ;
  \item le \emph{masque de pruning} doit être figé avant le fine-tuning, sans quoi la sparsité est détruite ;
  \item l'\emph{export ONNX} avec données externes peut paradoxalement augmenter la taille totale si le format n'est pas explicité.
\end{enumerate}

% --- 5.4 Cross-dataset ----------------------------------------------
\section{Cross-dataset : un plateau intrinsèque}
\label{sec:disc:cross}

La perte cross-dataset observée n'est pas anecdotique. Elle s'explique principalement par deux facteurs :
\begin{itemize}
  \item la \textbf{distribution des durées d'épisodes} est différente entre AFDB (épisodes courts fréquents) et LTAFDB (épisodes longs et soutenus) ;
  \item la \textbf{fréquence d'échantillonnage} est différente (250~Hz contre 128~Hz), ce qui modifie la résolution des intervalles \RR{}.
\end{itemize}

L'écart entre une cible métier ambitieuse (\(F_1 \geq 0{,}95\)) et le plafond observé (\(F_1 \approx 0{,}75\text{--}0{,}80\)) traduit la limite de l'information portée par les \RR{} seuls : sans la morphologie de l'\ECG{}, certains épisodes restent indissociables d'un rythme sinusal irrégulier.

% --- 5.5 Limitations ------------------------------------------------
\section{Limitations}
\label{sec:disc:limites}

Le présent travail comporte plusieurs limitations.

\textbf{Données.} L'effectif d'AFDB est restreint (25~patients) et n'est pas représentatif de la diversité des conditions cliniques rencontrées en pratique.

\textbf{Annotations.} Les positions d'ondes R ont été extraites des annotations expertes ; un détecteur de R automatique introduirait un biais supplémentaire en conditions réelles.

\textbf{Information utilisée.} L'utilisation des seuls intervalles \RR{}, choisie pour sa simplicité de déploiement, exclut par construction la morphologie de l'onde P.

\textbf{Évaluation.} L'évaluation s'est concentrée sur deux jeux de données ; un transfert vers des bases supplémentaires (CPSC2018, Icentia11k, etc.) reste à effectuer.

% =====================================================================
%  Chapitre 6 — Conclusion et perspectives
%  Règles (Sayadi 5 / Buttler) :
%   - récapituler brièvement le cheminement et les conclusions intermédiaires
%   - énumérer les propositions
%   - terminer par limitations + perspectives futures
% =====================================================================
\chapter{Conclusion et perspectives}
\label{ch:conclusion}

% --- 6.1 Récapitulatif ----------------------------------------------
\section{Récapitulatif du travail réalisé}
\label{sec:concl:recap}

Ce mémoire a proposé une chaîne complète de détection de la fibrillation auriculaire à partir des seuls intervalles \RR{}, structurée en six phases.

La \textbf{phase~1} a établi un protocole d'évaluation patient-level rigoureux et caractérisé la forte hétérogénéité inter-patients de la base AFDB.

La \textbf{phase~2} a comparé quatre méthodes de référence (règle métier, forêt aléatoire, gradient boosting, \CNN{} 1D) sous ce même protocole.

La \textbf{phase~3} a conçu une architecture \CNN--\LSTM{} optimisée par recherche bayésienne, accompagnée d'une étude d'ablation.

La \textbf{phase~4} a évalué huit variantes de compression et identifié la combinaison optimale (quantification statique INT8) qui atteint 87~KB d'empreinte mémoire sans dégradation significative.

La \textbf{phase~5} a quantifié la perte cross-dataset par transfert AFDB~$\rightarrow$~LTAFDB.

La \textbf{phase~6} a livré un tableau de bord interactif \textit{Streamlit} pour l'exploration patient par patient des décisions du modèle.

% --- 6.2 Apports principaux -----------------------------------------
\section{Apports principaux}
\label{sec:concl:apports}

Les apports de ce travail sont les suivants :
\begin{itemize}[leftmargin=*]
  \item la mise en évidence d'un \keyword{plateau intrinsèque} de \(F_1 \approx 0{,}75\text{--}0{,}80\) pour la classification patient à partir des seuls \RR{}, indépendant de la famille de classifieur ;
  \item une étude de compression aboutie permettant un déploiement embarqué (modèle de 87~KB) ;
  \item une mesure honnête de la robustesse cross-dataset, peu documentée jusqu'ici ;
  \item un outil d'exploration interactive favorisant la transparence et la pédagogie auprès d'un public clinicien.
\end{itemize}

% --- 6.3 Limitations ------------------------------------------------
\section{Limitations}
\label{sec:concl:limites}

Les principales limitations du travail, déjà discutées au chapitre~\ref{ch:discussion}, sont rappelées ici :
\begin{itemize}
  \item un effectif patient restreint sur AFDB ;
  \item une dépendance aux annotations expertes pour les positions de R ;
  \item l'utilisation exclusive du signal \RR{}, qui ignore la morphologie ;
  \item un nombre limité de jeux de données externes pour l'évaluation de la robustesse.
\end{itemize}

% --- 6.4 Perspectives -----------------------------------------------
\section{Perspectives}
\label{sec:concl:perspectives}

Plusieurs extensions de ce travail sont envisagées.

\textbf{Court terme.}
\begin{itemize}
  \item Réaliser un déploiement de l'application sur Streamlit Community Cloud pour rendre la démonstration accessible publiquement.
  \item Compléter les ré-exports des figures en \texttt{.pdf} 300 dpi pour la qualité d'impression du présent mémoire.
\end{itemize}

\textbf{Moyen terme.}
\begin{itemize}
  \item Ajouter une seconde branche d'entrée portant la morphologie de l'\ECG{} (onde P, complexe QRS), afin de tester l'hypothèse du plateau intrinsèque.
  \item Évaluer le modèle final sur d'autres jeux de données publics (CPSC2018, Icentia11k) pour étendre la mesure de robustesse.
\end{itemize}

\textbf{Long terme.}
\begin{itemize}
  \item Industrialiser le pipeline avec un test prospectif sur un dispositif embarqué (microcontrôleur ARM Cortex-M, par exemple).
  \item Étudier une approche d'\emph{adaptation de domaine} (fine-tuning court par patient ou par dataset cible) pour combler l'écart cross-dataset.
\end{itemize}


% ----- Bibliographie ---------------------------------------------------
\cleardoublepage
\phantomsection
\addcontentsline{toc}{chapter}{Références bibliographiques}
\printbibliography[title={Références bibliographiques}]

% ----- Annexes ---------------------------------------------------------
\appendix
% =====================================================================
%  Annexes (Buttler) : éléments externes au rapport, lecture optionnelle.
% =====================================================================

\chapter{Liste des hyperparamètres optimisés}
\label{ann:hyperparams}

Le tableau suivant détaille l'espace de recherche utilisé par Optuna en phase~3.

\begin{table}[H]
  \centering
  \caption{Espace de recherche Optuna.}
  \label{tab:ann:optuna}
  \begin{tabular}{lll}
    \toprule
    \textbf{Hyperparamètre} & \textbf{Plage} & \textbf{Type} \\
    \midrule
    Taille de fenêtre \(W\)        & [30, 60, 90, 120] & catégoriel \\
    Filtres conv. couche 1         & [16, 32, 64]      & catégoriel \\
    Filtres conv. couche 2         & [32, 64, 128]     & catégoriel \\
    Unités LSTM                    & [32, 64, 128]     & catégoriel \\
    Dropout                        & [0,1 ; 0,5]       & continu (log) \\
    Taux d'apprentissage           & [1e-5 ; 1e-3]     & continu (log) \\
    Taille de \textit{batch}       & [32, 64, 128]     & catégoriel \\
    \bottomrule
  \end{tabular}
\end{table}

\chapter{Reproductibilité et code source}
\label{ann:reproductibilite}

Le code complet du projet est disponible publiquement à l'adresse suivante :

\medskip
\centerline{\url{[À COMPLÉTER : URL du dépôt GitHub]}}
\medskip

Le fichier \texttt{requirements-lock.txt} fixe les versions exactes des bibliothèques utilisées. La commande \texttt{make demo} lance l'application interactive en local après installation des dépendances.

\chapter{Captures d'écran complémentaires}
\label{ann:screenshots}

\begin{figure}[H]
  \centering
  \figplaceholder[0.8\linewidth][5cm]{Capture — onglet Model Showdown}
  \caption{Onglet \emph{Model Showdown} de l'application interactive.}
\end{figure}

\begin{figure}[H]
  \centering
  \figplaceholder[0.8\linewidth][5cm]{Capture — onglet Compression Lab}
  \caption{Onglet \emph{Compression Lab}.}
\end{figure}

\begin{figure}[H]
  \centering
  \figplaceholder[0.8\linewidth][5cm]{Capture — onglet What-If}
  \caption{Onglet \emph{What-If}.}
\end{figure}


\end{document}
